Họ tên: bàn an bình 
Mã SV: 2071020026
Lớp: CNTT 20-04

1. KẾT QUẢ KIỂM THỬ T01–T08

T01 – Kiểm tra trang chủ
- Kỳ vọng: Trang Portfolio hiển thị đầy đủ nội dung.
- Thực tế: Trang chủ mở được và hiển thị các phần nội dung.
- Kết quả: Đạt.
- Minh chứng: Ảnh chụp trang chủ.

T02 – Kiểm tra menu điều hướng
- Kỳ vọng: Các liên kết trên menu chuyển đến đúng mục.
- Thực tế: Các liên kết Giới thiệu, Kỹ năng, Dự án, Lộ trình và Liên hệ hoạt động đúng.
- Kết quả: Đạt.
- Minh chứng: Ảnh chụp menu điều hướng.

T03 – Kiểm tra giao diện CSS
- Kỳ vọng: CSS được áp dụng và bố cục hiển thị đúng.
- Thực tế: Giao diện hiển thị bình thường, không bị mất định dạng.
- Kết quả: Đạt.
- Minh chứng: Ảnh chụp giao diện Portfolio.

T04 – Kiểm tra hiệu ứng hover của nút
- Kỳ vọng: Nút thay đổi giao diện khi rê chuột và có phản hồi khi nhấn.
- Thực tế: Nút có hiệu ứng tương tác theo CSS.
- Kết quả: Đạt.
- Minh chứng: Ảnh chụp nút khi rê chuột.

T05 – Kiểm tra hiệu ứng thẻ dự án
- Kỳ vọng: Thẻ dự án có hiệu ứng khi rê chuột hoặc focus vào thành phần bên trong.
- Thực tế: Thẻ dự án có hiệu ứng tương tác.
- Kết quả: Đạt.
- Minh chứng: Ảnh chụp thẻ dự án.

T06 – Kiểm tra điều hướng bằng bàn phím
- Kỳ vọng: Có thể dùng phím Tab để di chuyển qua các liên kết và nút.
- Thực tế: Các phần tử tương tác có thể được chọn bằng bàn phím và có chỉ báo focus.
- Kết quả: Đạt.
- Minh chứng: Ảnh chụp phần tử đang focus.

T07 – Kiểm tra dữ liệu form không hợp lệ
- Kỳ vọng: Form không được gửi khi bỏ trống trường bắt buộc hoặc nhập dữ liệu không hợp lệ.
- Thực tế: Trình duyệt yêu cầu nhập các trường bắt buộc và kiểm tra định dạng email.
- Kết quả: Đạt.
- Minh chứng: Ảnh chụp thông báo kiểm tra form.

T08 – Kiểm tra dữ liệu form hợp lệ
- Kỳ vọng: Form chuyển đến trang kết quả khi dữ liệu đáp ứng điều kiện.
- Thực tế: Form chuyển đến ket-qua.html khi dữ liệu hợp lệ. Đây là form mô phỏng, không gửi email hoặc lưu dữ liệu vào cơ sở dữ liệu.
- Kết quả: Đạt.
- Minh chứng: Ảnh chụp trang kết quả.

2. GHI NHẬN MỘT LỖI VÀ CÁCH KHẮC PHỤC

Bước tái hiện: Mở index.html trên trình duyệt và kiểm tra giao diện khi CSS không được áp dụng.

Nguyên nhân: Đường dẫn đến tệp style.css không khớp với cấu trúc thư mục hoặc tệp CSS chưa được lưu.

Cách sửa: Kiểm tra đường dẫn trong thẻ link của index.html. Nếu tệp CSS nằm trong thư mục css thì sử dụng `href="css/style.css"`. Lưu lại các tệp và tải lại trang.

Kiểm tra lại: Sau khi kiểm tra đường dẫn và tải lại trang, giao diện hiển thị đúng định dạng CSS.

3. ÁP DỤNG SELECTOR CSS VÀO PORTFOLIO

Selector mẫu: a:focus-visible

Selector áp dụng: .menu a:focus-visible

Lý do: Selector này áp dụng chỉ báo focus cho các liên kết trong menu khi người dùng điều hướng bằng bàn phím. Nhờ đó, người dùng dễ nhận biết liên kết đang được chọn và sử dụng trang thuận tiện hơn.

4. SỬ DỤNG AI VÀ TỰ KIỂM CHỨNG

Phần sử dụng AI: Tham khảo AI để hỗ trợ điền các thuộc tính HTML5 của biểu mẫu, định dạng mã HTML và kiểm tra các selector CSS.

Cách tự kiểm chứng: Đọc lại mã nguồn, mở Portfolio trên trình duyệt, kiểm tra menu điều hướng, hiệu ứng hover, focus bằng bàn phím và thử form với dữ liệu không hợp lệ cũng như dữ liệu hợp lệ. Đối chiếu kết quả hiển thị thực tế và chụp ảnh minh chứng trước khi nộp bài.